\section{Scales, density tails, and ground-state paths}\label{sec:setup}

We first choose units in which the interaction range is small, the particle
density is large, and the healing length is of order one. Two local
density-tail estimates will then supply the moment and pair bounds used
throughout the paper. Their first use is to control the boundary correction
in the energy calculation of Section~\ref{sec:energy}.

\subsection{Ground states and scaled variables}

For finite $N,L$, the potential in the Hamiltonian is a bounded
multiplication operator. The heat kernel on the labelled configuration
torus is strictly positive, and the Feynman--Kac formula preserves this
positivity. Compact resolvent and positivity improvement therefore give a
unique normalized positive ground function. Permutations of coordinates
and simultaneous translations commute with the Hamiltonian. Uniqueness
implies that this ground function is symmetric and translation invariant.
In particular the bosonic ground space is one-dimensional. We shall use
positivity when comparing spatial arithmetic means with root mean squares
in Section~\ref{sec:assembly}.

Let $R_v>0$ contain the support of the physical potential $v$, and use a
length unit $\sigma>0$. In the resulting coordinates set
\begin{equation}\label{eq:scaled-variables}
 \begin{gathered}
 K=L/\sigma,\qquad V=K^3,\qquad D=N/V,\qquad a=a_v/\sigma,\\
 \alpha=Da,\qquad r=R_v/\sigma,\qquad W(y)=\sigma^2v(\sigma y).
 \end{gathered}
\end{equation}
The scaled Hamiltonian has kinetic energy $-\Delta$ and pair potential
$W$, periodized on the torus of side $K$. Its normalized positive ground
function is denoted by $\Phi$, and
\[
 \mu(\dd X)=\Phi(X)^2\dd X.
\]
Integrals in all particle coordinates use Lebesgue measure. We write
$\langle O\rangle_\psi=\langle\psi,O\psi\rangle$; this notation does not
normalize a vector that is not already normalized.

Fix a sufficiently large constant $B_0$, depending on $v$, and put
\begin{equation}\label{eq:physical-base-scale}
 \sigma_*=B_0\rho^{-1/2},\qquad
 D_*=\rho\sigma_*^3=B_0^3\rho^{-1/2},\qquad
 \alpha_0=a_vB_0^2.
\end{equation}
Along the thermodynamic sequence we choose $\sigma\to\sigma_*$, as
specified below. Then $D\to D_*$ and $\alpha\to\alpha_0$. All estimates
are uniform for $\alpha$ in a fixed compact positive neighborhood of
$\alpha_0$. Constants may depend on the fixed potential and this
neighborhood. In these units,
\begin{equation}\label{eq:potential-scales}
 a\asymp r\asymp D^{-1},\qquad
 \int_{\R^3}W\le C D^{-1},\qquad 0\le W\le C D^2.
\end{equation}
A bounded measurable radial potential can be replaced by a Borel version
with its range and radial bounds holding everywhere. Such null changes
affect neither the operator nor the Brownian path integrals.

Choose small exponents $u>0$ and $0<w<u/10000$, and let
\begin{equation}\label{eq:partition-scales}
 R=\lceil D_*^w\rceil,\qquad
 B=R\lceil D_*^u/R\rceil,\qquad d=1/j\quad(j\in\N).
\end{equation}
The fixed microcell size $d$ resolves variation within cells.  The
slowly growing scale $R$ is used for the distant deletion comparison,
while the still larger $B$-boxes give the energy calculation room to
resolve variation on the $R$-scale.  For each fixed $j$, the dilute
hierarchy is $r\ll d\le1\ll R\ll B$.  All partition sizes are fixed
during the thermodynamic limit, when $K\to\infty$; $j$ is sent to
infinity only after the estimates at each fixed $j$.

The restrictions on $u$ arise in the energy estimates; it may be taken
smaller than $10^{-9}$ and decreased further there. The path exponent
$x>0$ will be chosen after $u,w$. At fixed $\rho$, choose $K$ to be the
nearest positive multiple of $B$ to $L/\sigma_*$ and set $\sigma=L/K$.
This gives the asserted convergence of $\sigma$ on a tail of every
thermodynamic sequence.

We use nested cubic partitions $\cC_l$ of the scaled torus for
$l=B,R,1,d$. Unit cells $U_z$ are centered at
$z\in\cL=(\Z/K\Z)^3$, with compatible half-open conventions on
boundaries. If $A$ is a spatial set, $n_A$ denotes its particle count.
Let $P_l$ be the one-body orthogonal projection onto functions constant
on every cell of $\cC_l$, and set
\begin{equation}\label{eq:cell-projections}
 W_l=1-P_l,\qquad \dGamma(O)=\sum_{i=1}^N O_i.
\end{equation}
Thus $W_l$ is a projection, whereas $W$ without a subscript is the pair
potential. Averages over cells or cell pairs are uniform unless a
different measure is specified.

\paragraph{Limit convention.}
We always take the thermodynamic limsup at fixed $\rho$ first, followed
by $\rho\downarrow0$, equivalently $D_*\to\infty$. A term $o_K(1)$
tends to zero in the first limit; a dilute $o(1)$ tends to zero after
that limsup. Power bounds in $D$ hold for all sufficiently large $D_*$
and sufficiently far in the corresponding inner sequence. Terms tending
to zero only in the inner limit can therefore be absorbed into a stated
power bound on a tail. No uniform rate in that sequence is required.
Constants and thresholds in estimates involving $d$ may depend on the
fixed integer $j$.

\subsection{Local density tails}

We use two results from the bounded-interaction companion~\cite{C}.
We record both the physical form of the input and the consequence needed
here, so that the dependence on temperature and volume is explicit.
For a canonical Gibbs state let its position law be the diagonal of its
density matrix.

\begin{proposition}[Density-tail input from \cite{C}]\label{prop:density-input}
Let $v$ satisfy the hypotheses of Theorem~\ref{thm:main}. There is a
sufficiently large fixed $B_0$ such that, for all sufficiently small
$\rho>0$, the following statements hold at physical base length
$b=B_0\rho^{-1/2}$ and count scale $n=\rho b^3$.
\begin{enumerate}[label=\textup{(\roman*)}]
\item If the particle density is at most $2\rho$, let
$A=\sum_i g(x_i)$, where $0\le g\le1$, the support of $g$ is covered
by $J$ base cubes, and $|\nabla g|^2\le C_gb^{-2}g$ almost everywhere
for a prescribed fixed bound on $C_g$. Then
\[
 \Pp(A>M)\le C e^{-cM},\qquad M\ge CnJ.
\]
Sharp counts have such majorants after a fixed enlargement and a fixed
increase in $J$.
\item If the density lies in $[\rho/2,2\rho]$ and $Q_1,\ldots,Q_J$
are cubes of side $b$ with centers separated by a sufficiently large
fixed multiple of $b$, then
\[
 \Pp\left(\sum_{i=1}^J n_{Q_i}\le\epsilon_0 nJ\right)
 \le e^{-c_1 nJ}
\]
for fixed $\epsilon_0,c_1>0$.
\end{enumerate}
Both assertions hold in all sufficiently large volumes, uniformly down
to arbitrarily small positive temperatures. Their volume thresholds do
not grow as temperature decreases. Fixed positive multiplicative changes
in cube sides are allowed after changing constants, and $B_0$ can be
increased to meet any finite list of lower bounds.
\end{proposition}

The density-tail inputs are supplied by the localized density and occupation
estimates of~\cite[Lemmas~4.2 and~4.3]{C}; for the upper tail we take
exponential parameter one.
Their proofs use local energy and entropy estimates and a localized
replacement argument. They do not require a condensation or depletion
conclusion. Their temperature uniformity lets us apply them to the
ground state at each fixed finite volume: canonical Gibbs states converge
in trace norm to its rank-one projection as temperature tends to zero.

\begin{lemma}[Scaled count tails]\label{lem:count-tails}
After fixing $B_0$ and decreasing $\rho$ if necessary, there are positive
constants $c,c',C$ such that, for every nonempty $A\subset\cL$,
\begin{equation}\label{eq:count-tails}
 \begin{aligned}
 \Pp_\mu\left(\sum_{z\in A}n_{U_z}\ge b\right)
   &\le C e^{-cb},&& b\ge CD|A|,\\
 \Pp_\mu\left(n_{U_z}<c'D\text{ for every }z\in A\right)
   &\le e^{-cD|A|}.
 \end{aligned}
\end{equation}
\end{lemma}
\begin{proof}
The physical side of a scaled unit cell is $\sigma\to b$, and
$D\asymp D_*$. A fixed enlargement covers any union of the given unit
cells by $O(|A|)$ base cubes, proving the first bound from
Proposition~\ref{prop:density-input}(i). For the second, a bounded coloring
of the unit lattice extracts a subset of $A$ of size at least $c_0|A|$
whose cells have the separation required in part (ii). If every original
count is below $c'D$, their sum on this subset is below the threshold of
part (ii) when $c'$ is small enough. The allowed fixed side adjustments
handle either sign of $\sigma-b$. The actual physical density eventually
lies in $[\rho/2,2\rho]$. This proves both claims uniformly on a tail
of the thermodynamic sequence.
\end{proof}

For $A\subset\cL$ define
\[
 F_A(X)=\sum_{i=1}^N e^{-d_K(x_i,A)},\qquad F_\varnothing=0,
\]
where $d_K$ denotes torus distance and the sites of $A$ are regarded as
points of the torus. The tails imply the exponential moment bound
\begin{equation}\label{eq:count-mgf}
 \log\E_\mu e^{hF_A}\le ChD|A|\qquad(0\le h\le h_*),
\end{equation}
for a fixed $h_*>0$. Indeed the upper tail first gives this statement
for the sharp count in a site set and one small positive value of $h$;
convexity extends it to the interval from zero. The sites at distance at
most $m+1$ from $A$ number at most $C(1+m)^3|A|$. Write the exponentially
weighted count as a sum of counts in these enlargements, with exponentially
decaying coefficients. Jensen's inequality with normalized summable weights
and the sharp-count bound prove~\eqref{eq:count-mgf}. The same proof works
if distance is multiplied by any fixed positive constant.

In particular, each fixed $q$th moment of a count in a fixed number of
unit cells is $O(D^q)$. Counts in polynomially many cells have polynomial
fixed moments. If one of those cells exceeds a sufficiently large fixed
multiple of $D$, its probability is exponentially small in $D$.
Cauchy--Schwarz therefore makes every fixed product of these counts on
that exceptional event smaller than every power of $D^{-1}$. We will use
this observation when removing count cutoffs.

\subsection{Ground-state paths and local pair bounds}

Let $E_m$ be the ground energy of $m$ particles on the scaled torus,
with the same potential $W$. Adjoining a constant last coordinate to a
normalized $(m-1)$-particle trial function gives
\begin{equation}\label{eq:removal-energy}
 E_m-E_{m-1}\le\frac{m-1}{V}\int W\le C
 \quad(m\le N),\qquad E_N-E_{N-m}\le Cm.
\end{equation}
This trial need not be symmetric: the positive ground energy on the
labelled space equals the bosonic ground energy by the uniqueness already
proved.

For $0<T\le1$, consider independent torus Brownian motions on $[-T,T]$
with generator $\Delta$ and initial Lebesgue measure. Weight their joint
law by
\begin{equation}\label{eq:ground-path-law}
 e^{2TE_N}\Phi(\omega(-T))\Phi(\omega(T))
 \prod_{i<k}\exp\left[-\int_{-T}^T
              W(\omega_i(q)-\omega_k(q))\dd q\right].
\end{equation}
Here and below the pair potential in torus expressions is periodized.
The Feynman--Kac formula makes~\eqref{eq:ground-path-law} a probability
measure with marginal $\mu$ at every time. Restrictions to shorter slabs
are consistent, by the ground-state eigenfunction identity.

\begin{lemma}[Free-path domination]\label{lem:path-domination}
If $m$ specified labels carry nonnegative single-path weights $H_i$ that
are invariant under time reversal, then under~\eqref{eq:ground-path-law},
\begin{equation}\label{eq:path-domination}
 \E\prod_{i=1}^m H_i(\omega_i)
 \le e^{Cm}\int\mu(\dd Y)
                  \prod_{i=1}^m\E^{Y_i}_{\mathrm{free}}H_i.
\end{equation}
A weight without reversal symmetry can be replaced by the sum of itself
and its reversal.
\end{lemma}
\begin{proof}
Delete every interaction incident to the specified labels; nonnegativity
of $W$ increases the integral. Fix their initial and final coordinates
$y,z$. The remaining semigroup has norm $e^{-2TE_{N-m}}$, so the integral
in the other labels is bounded by that norm times $g(y)g(z)$, where
$g(y)$ is the $L^2$ norm of the slice $\Phi(y,\cdot)$. The product of
the free weighted endpoint kernels is symmetric in $y,z$ by reversal
invariance. Thus $g(y)g(z)\le(g(y)^2+g(z)^2)/2$ replaces the slice
product by its initial square after integration. Since $g(y)^2$ is the
specified-coordinate marginal of $\mu$, the result follows from
\eqref{eq:removal-energy}. Truncation proves the assertion for unbounded
nonnegative weights.
\end{proof}

Apply this lemma on the longer slab $[-1,1]$. At any time $|q|\le1/2$,
the probability for a free path started at $Y_i$ to lie in a cube of
side $h\le1$ within a fixed distance of a site $z$ is bounded by
$Ch^3e^{-2d_K(Y_i,z)}$, by the torus heat kernel. The bound also holds
for the reversed observation. Summing~\eqref{eq:path-domination} over
distinct labels and using~\eqref{eq:count-mgf} shows that two specified
such cubes have expected ordered pair count at most $CD^2h^6$.
Covering a fixed site neighborhood by $h$-cubes and summing over the
bounded number of neighboring cubes for each one gives
\begin{equation}\label{eq:pair-bound}
 \E\#\{(i,k):i\ne k,\ X_i,X_k\text{ in the neighborhood},\
                    d_K(X_i,X_k)\le h\}\le CD^2h^3.
\end{equation}
These bounds are uniform in the site and in the observation time on the
central slab. They provide the required short-distance control using only
the imported count tails and the elementary removal-energy estimate.
